\begin{bt}%[9D6H18-2]
% nguon: Đề tham khảo Tuyển sinh 10 TP.HCM 2025-2026 – Quận 4 – Đề 1 – Bài 1
Cho parabol $(P): y = -\dfrac{1}{2}x^2$.
\begin{enumerate}[a)]
	\item Vẽ đồ thị $(P)$ trên hệ trục toạ độ $Oxy$.
	\item Tìm toạ độ những điểm $A$ thuộc $(P)$ biết rằng tung độ của nó gấp đôi hoành độ.
\end{enumerate}
\loigiai{
\begin{enumerate}[a)]
	\item Bảng giá trị tương ứng giữa $x$ và $y$ của hàm số $y = -\dfrac{1}{2}x^2$:
	\begin{center}
	\begin{tabular}{|c|c|c|c|c|c|}
	\hline
	$x$ & $-4$ & $-2$ & $0$ & $2$ & $4$ \\
	\hline
	$y = -\dfrac{1}{2}x^2$ & $-8$ & $-2$ & $0$ & $-2$ & $-8$ \\
	\hline
	\end{tabular}
	\end{center}
	Đồ thị $(P)$ là một parabol có đỉnh tại gốc toạ độ $O(0; 0)$, nhận trục tung $Oy$ làm trục đối xứng và nằm phía dưới trục hoành.
	\item Điểm $A(x; y)$ thuộc $(P)$ có tung độ gấp đôi hoành độ thoả mãn:
	$y = 2x \Leftrightarrow -\dfrac{1}{2}x^2 = 2x \Leftrightarrow x^2 + 4x = 0 \Leftrightarrow x(x + 4) = 0 \Leftrightarrow \left[\begin{aligned} x &= 0 \Rightarrow y = 0 \\ x &= -4 \Rightarrow y = -8. \end{aligned}\right.$\\
	Vậy có hai điểm thoả mãn là $A_1(0; 0)$ và $A_2(-4; -8)$.
\end{enumerate}
}
\end{bt}

\begin{bt}%[9D6H20-2]
% nguon: Đề tham khảo Tuyển sinh 10 TP.HCM 2025-2026 – Quận 4 – Đề 1 – Bài 2
Cho phương trình $3x^2 - 7x + 3 = 0$.
\begin{enumerate}[a)]
	\item Chứng minh phương trình trên có hai nghiệm phân biệt.
	\item Không giải phương trình, gọi $x_1, x_2$ là hai nghiệm của phương trình đã cho, hãy tính giá trị của biểu thức:
	$A = \dfrac{18x_2}{5x_1 - x_2 + \dfrac{7}{3}} + \dfrac{18x_1}{5x_2 - x_1 + \dfrac{7}{3}}$.
\end{enumerate}
\loigiai{
\begin{enumerate}[a)]
	\item Xét phương trình $3x^2 - 7x + 3 = 0$, ta có:
	$\Delta = b^2 - 4ac = (-7)^2 - 4\cdot 3\cdot 3 = 49 - 36 = 13 > 0$.\\
	Vì $\Delta > 0$ nên phương trình luôn có hai nghiệm phân biệt $x_1, x_2$.
	\item Theo định lí Viète, ta có:
	$\begin{cases} x_1 + x_2 = -\dfrac{b}{a} = \dfrac{7}{3} \\ x_1 x_2 = \dfrac{c}{a} = \dfrac{3}{3} = 1. \end{cases}$\\
	Thay $\dfrac{7}{3} = x_1 + x_2$ vào mẫu thức của biểu thức $A$:
	\begin{itemize}
		\item Mẫu thức thứ nhất: $5x_1 - x_2 + \dfrac{7}{3} = 5x_1 - x_2 + (x_1 + x_2) = 6x_1$.
		\item Mẫu thức thứ hai: $5x_2 - x_1 + \dfrac{7}{3} = 5x_2 - x_1 + (x_1 + x_2) = 6x_2$.
	\end{itemize}
	Khi đó biểu thức $A$ trở thành:
	$A = \dfrac{18x_2}{6x_1} + \dfrac{18x_1}{6x_2} = 3\left(\dfrac{x_2}{x_1} + \dfrac{x_1}{x_2}\right) = 3\cdot\dfrac{x_1^2 + x_2^2}{x_1 x_2} = 3\cdot\dfrac{(x_1 + x_2)^2 - 2x_1 x_2}{x_1 x_2}$.\\
	Thay $x_1 + x_2 = \dfrac{7}{3}$ và $x_1 x_2 = 1$ vào, ta được:
	$A = 3\cdot\dfrac{\left(\dfrac{7}{3}\right)^2 - 2\cdot 1}{1} = 3\left(\dfrac{49}{9} - 2\right) = 3\cdot\dfrac{31}{9} = \dfrac{31}{3}$.\\
	Vậy $A = \dfrac{31}{3}$.
\end{enumerate}
}
\end{bt}

\begin{bt}%[9D8H26-3]
% nguon: Đề tham khảo Tuyển sinh 10 TP.HCM 2025-2026 – Quận 4 – Đề 1 – Bài 3
Bảng xếp hạng bóng đá bảng B gồm $5$ đội bóng thi đấu vòng tròn một lượt với kết quả số trận thắng, hoà, thua như sau:
\begin{itemize}
	\item Việt Nam: $4$ trận ($3$ thắng, $1$ hoà, $0$ thua), đạt $10$ điểm.
	\item Philippines: $4$ trận ($2$ thắng, $1$ hoà, $1$ thua), đạt $7$ điểm.
	\item Indonesia: $4$ trận ($1$ thắng, $1$ hoà, $2$ thua), đạt $4$ điểm.
	\item Myanmar: $4$ trận ($1$ thắng, $0$ hoà, $3$ thua), đạt $3$ điểm.
	\item Lào: $4$ trận ($0$ thắng, $0$ hoà, $4$ thua), đạt $0$ điểm.
\end{itemize}
Chọn ngẫu nhiên một đội bóng trong số $5$ đội bóng trên:
\begin{enumerate}[a)]
	\item Tính xác suất để đội bóng được chọn có đúng $1$ trận thắng.
	\item Tính xác suất để đội bóng được chọn có số điểm từ $4$ điểm trở lên.
\end{enumerate}
\loigiai{
Số phần tử của không gian mẫu khi chọn ngẫu nhiên $1$ đội bóng từ $5$ đội bóng là $n(\Omega) = 5$.
\begin{enumerate}[a)]
	\item Gọi biến cố $E$: ``Đội bóng được chọn có đúng $1$ trận thắng''.\\
	Trong bảng xếp hạng, các đội có đúng $1$ trận thắng là: Indonesia và Myanmar (gồm $2$ đội bóng).\\
	Số kết quả thuận lợi cho biến cố $E$ là $n(E) = 2$.\\
	Xác suất cần tìm là: $P(E) = \dfrac{n(E)}{n(\Omega)} = \dfrac{2}{5} = 0{,}4$.
	\item Gọi biến cố $F$: ``Đội bóng được chọn có số điểm từ $4$ điểm trở lên''.\\
	Các đội đạt từ $4$ điểm trở lên gồm: Việt Nam ($10$ điểm), Philippines ($7$ điểm) và Indonesia ($4$ điểm) (gồm $3$ đội bóng).\\
	Số kết quả thuận lợi cho biến cố $F$ là $n(F) = 3$.\\
	Xác suất cần tìm là: $P(F) = \dfrac{n(F)}{n(\Omega)} = \dfrac{3}{5} = 0{,}6$.
\end{enumerate}
}
\end{bt}

\begin{bt}%[9D1H2-1]
% nguon: Đề tham khảo Tuyển sinh 10 TP.HCM 2025-2026 – Quận 4 – Đề 1 – Bài 4
Một câu lạc bộ bóng rổ tổ chức tuyển chọn thành viên mới với quy định: Mỗi học sinh thực hiện $25$ lượt ném bóng. Mỗi lượt ném bóng vào rổ được cộng $5$ điểm, mỗi lượt ném không vào rổ bị trừ $2$ điểm. Ban đầu mỗi bạn có sẵn $50$ điểm. Để được chọn vào đội thì sau $25$ lượt ném, học sinh phải đạt tổng số điểm từ $100$ điểm trở lên.
\begin{enumerate}[a)]
	\item Một bạn học sinh thực hiện $25$ lượt ném gồm $18$ lượt bóng vào rổ và $7$ lượt bóng không vào rổ. Tính tổng số điểm của học sinh đó.
	\item Gọi $x$ là số lượt ném bóng vào rổ trong $25$ lượt ($x \in \mathbb{N}, 0 \le x \le 25$). Hãy biểu diễn tổng số điểm của học sinh theo $x$. Hỏi để được chọn vào đội, học sinh cần ném bóng vào rổ ít nhất bao nhiêu lượt?
\end{enumerate}
\loigiai{
\begin{enumerate}[a)]
	\item Tổng số điểm của học sinh sau khi thực hiện $18$ lượt vào rổ và $7$ lượt không vào rổ là:
	$50 + 18 \cdot 5 - 7 \cdot 2 = 50 + 90 - 14 = 126\text{ (điểm)}$.
	\item Trong $25$ lượt ném, nếu có $x$ lượt bóng vào rổ thì số lượt bóng không vào rổ là $25 - x$.\\
	Biểu thức biểu diễn tổng số điểm của học sinh theo $x$ là:
	$T(x) = 50 + 5x - 2(25 - x) = 50 + 5x - 50 + 2x = 7x\text{ (điểm)}$.\\
	Để học sinh được chọn vào đội, tổng số điểm phải đạt từ $100$ điểm trở lên:
	$7x \ge 100 \Leftrightarrow x \ge \dfrac{100}{7} \approx 14{,}28$.\\
	Vì $x$ là số nguyên dương và $x \le 25$ nên $x \ge 15$.\\
	Vậy học sinh cần thực hiện ít nhất $15$ lượt ném bóng vào rổ thì mới được chọn vào đội.
\end{enumerate}
}
\end{bt}

\begin{bt}%[9H0V31-3]
% nguon: Đề tham khảo Tuyển sinh 10 TP.HCM 2025-2026 – Quận 4 – Đề 1 – Bài 5
Một ống đựng các viên vitamin C có dạng hình trụ với đường kính đáy là $4\text{ cm}$ và chiều cao là $12\text{ cm}$. Mỗi viên vitamin C có dạng hình cầu với bán kính là $0{,}5\text{ cm}$.
\begin{enumerate}[a)]
	\item Tính thể tích của ống hình trụ và thể tích của một viên vitamin C (lấy $\pi \approx 3{,}14$, làm tròn kết quả đến hàng phần trăm của $\text{cm}^3$).
	\item Biết rằng khoảng không khí giữa các viên thuốc chiếm $10\%$ thể tích của ống. Tính số viên vitamin C tối đa có thể chứa được trong chiếc ống đó.
\end{enumerate}
\loigiai{
\begin{enumerate}[a)]
	\item Bán kính đáy của ống hình trụ là: $R = \dfrac{4}{2} = 2\text{ (cm)}$.\\
	Thể tích của ống hình trụ là:
	$V_{\text{ống}} = \pi R^2 h = 3{,}14 \cdot 2^2 \cdot 12 = 150{,}72\text{ (cm}^3\text{)}$.\\
	Thể tích của một viên vitamin C hình cầu bán kính $r = 0{,}5\text{ cm}$ là:
	$V_{\text{viên}} = \dfrac{4}{3}\pi r^3 = \dfrac{4}{3} \cdot 3{,}14 \cdot (0{,}5)^3 = \dfrac{4}{3} \cdot 3{,}14 \cdot 0{,}125 \approx 0{,}52\text{ (cm}^3\text{)}$.
	\item Vì lớp không khí chiếm $10\%$ thể tích ống nên thể tích thực tế dành cho các viên vitamin C là $90\%$ thể tích ống:\\
	$V_{\text{thuốc}} = 90\% \cdot V_{\text{ống}} = 0{,}9 \cdot 150{,}72 = 135{,}648\text{ (cm}^3\text{)}$.\\
	Số viên vitamin C tối đa có thể chứa trong ống là:
	$N = \dfrac{V_{\text{thuốc}}}{V_{\text{viên}}} = \dfrac{135{,}648}{\frac{4}{3}\cdot 3{,}14\cdot 0{,}125} \approx \dfrac{135{,}648}{0{,}5233} \approx 259{,}2$.\\
	Vì số viên thuốc là số nguyên nên ống có thể chứa tối đa $259$ viên vitamin C.
\end{enumerate}
}
\end{bt}

\begin{bt}%[9D1V3-4]
% nguon: Đề tham khảo Tuyển sinh 10 TP.HCM 2025-2026 – Quận 4 – Đề 1 – Bài 6
Cô Mai nhập $210\text{ bao gạo}$ gồm hai loại: bao gạo loại 1 (mỗi bao nặng $25\text{ kg}$, bán với giá $22\,500\text{ đồng/kg}$) và bao gạo loại 2 (mỗi bao nặng $50\text{ kg}$, bán với giá $18\,200\text{ đồng/kg}$) với tổng chi phí nhập hàng là $99\,700\,000\text{ đồng}$. Do thời tiết ẩm ướt, $5\%$ khối lượng gạo loại 1 và $6\%$ khối lượng gạo loại 2 bị hỏng không thể bán được. Tổng khối lượng gạo còn lại bán được là $6{,}15\text{ tấn} = 6150\text{ kg}$. Hỏi sau khi bán hết toàn bộ số gạo còn lại, cô Mai lãi hay lỗ bao nhiêu tiền?
\loigiai{
Gọi số bao gạo loại 1 và loại 2 cô Mai đã nhập lần lượt là $x$ và $y$ ($x, y \in \mathbb{N}^*, x, y < 210$).\\
Vì tổng số bao gạo nhập về là $210$ bao nên ta có phương trình:
$x + y = 210$.\\
Tổng khối lượng gạo ban đầu của loại 1 là $25x\text{ (kg)}$, sau khi hỏng $5\%$ thì khối lượng bán được là:
$25x \cdot (100\% - 5\%) = 25x \cdot 0{,}95 = 23{,}75x\text{ (kg)}$.\\
Tổng khối lượng gạo ban đầu của loại 2 là $50y\text{ (kg)}$, sau khi hỏng $6\%$ thì khối lượng bán được là:
$50y \cdot (100\% - 6\%) = 50y \cdot 0{,}94 = 47y\text{ (kg)}$.\\
Vì tổng khối lượng gạo còn lại bán được là $6150\text{ kg}$ nên ta có phương trình:
$23{,}75x + 47y = 6150$.\\
Ta có hệ phương trình:
$\begin{cases} x + y = 210 \\ 23{,}75x + 47y = 6150 \end{cases} \Leftrightarrow \begin{cases} y = 210 - x \\ 23{,}75x + 47(210 - x) = 6150 \end{cases} \Leftrightarrow \begin{cases} 23{,}25x = 3720 \\ y = 210 - x \end{cases} \Leftrightarrow \begin{cases} x = 160 \\ y = 50 \end{cases}$ (thoả mãn điều kiện).\\
Khối lượng gạo loại 1 bán được là: $160 \cdot 23{,}75 = 3800\text{ (kg)}$.\\
Doanh thu từ gạo loại 1 là: $3800 \cdot 22\,500 = 85\,500\,000\text{ (đồng)}$.\\
Khối lượng gạo loại 2 bán được là: $50 \cdot 47 = 2350\text{ (kg)}$.\\
Doanh thu từ gạo loại 2 là: $2350 \cdot 18\,200 = 42\,770\,000\text{ (đồng)}$.\\
Tổng doanh thu thu được từ việc bán gạo là:
$85\,500\,000 + 42\,770\,000 = 128\,270\,000\text{ (đồng)}$.\\
Số tiền lãi của cô Mai sau khi trừ chi phí nhập hàng là:
$128\,270\,000 - 99\,700\,000 = 28\,570\,000\text{ (đồng)}$.\\
Vậy cô Mai có lãi và số tiền lãi là $28\,570\,000\text{ đồng}$.
}
\end{bt}

\begin{bt}%[9H9C29-3]
% nguon: Đề tham khảo Tuyển sinh 10 TP.HCM 2025-2026 – Quận 4 – Đề 1 – Bài 7
Cho tam giác $ABC$ vuông tại $A$ có $AB = 8\text{ cm}$. Vẽ đường tròn tâm $O$ đường kính $AB$ cắt cạnh $BC$ tại điểm $D$ ($D \neq B$). Từ $A$ kẻ dây cung $AE$ của $(O)$ vuông góc với $OC$ tại $H$. Tiếp tuyến tại $D$ của đường tròn $(O)$ cắt $CA$ và $CE$ lần lượt tại $M$ và $I$.
\begin{enumerate}[a)]
	\item Chứng minh $4$ điểm $A, C, D, H$ cùng thuộc một đường tròn có tâm là $M$.
	\item Chứng minh $DH \perp DE$.
	\item Giả sử diện tích tứ giác $ACEO$ gấp $3$ lần diện tích tam giác $MOI$. Hãy tính diện tích tam giác $DHE$.
\end{enumerate}
\loigiai{
\begin{enumerate}[a)]
	\item Điểm $D$ thuộc đường tròn $(O)$ đường kính $AB$ nên $\widehat{ADB} = 90^\circ$, suy ra $AD \perp BC$ tại $D$.\\
	Do đó $\Delta ADC$ vuông tại $D$, tâm đường tròn ngoại tiếp $\Delta ADC$ là trung điểm $M$ của cạnh huyền $AC$, suy ra $MA = MC = MD$.\\
	Mặt khác, $OC \perp AE$ tại $H$ nên $H$ là trung điểm của dây $AE$.\\
	Xét tam giác $OAE$ cân tại $O$, có $OH \perp AE$ nên $OC$ là tia phân giác của $\widehat{AOE}$, suy ra $\Delta OAC = \Delta OEC$ (c.g.c).\\
	Do đó $\widehat{OEC} = \widehat{OAC} = 90^\circ$, suy ra $CE$ là tiếp tuyến của $(O)$.\\
	Tam giác $AHC$ vuông tại $H$ có trung tuyến $MH$ ứng với cạnh huyền $AC$ nên $MH = MA = MC$.\\
	Vậy bốn điểm $A, C, D, H$ cùng thuộc đường tròn tâm $M$ đường kính $AC$.
	\item Vì bốn điểm $A, C, D, H$ cùng thuộc đường tròn tâm $M$ nên tứ giác $ACDH$ nội tiếp đường tròn, suy ra $\widehat{CHD} = \widehat{CAD}$.\\
	Mặt khác, trong đường tròn $(O)$, góc nội tiếp $\widehat{CED}$ và góc tạo bởi tia tiếp tuyến và dây cung $\widehat{CAD}$ cùng chắn cung $AD$, nên $\widehat{CED} = \widehat{CAD}$.\\
	Từ đó suy ra $\widehat{CHD} = \widehat{CED}$, dẫn đến tứ giác $CDHE$ nội tiếp.\\
	Vì $\widehat{CDE} = 90^\circ$ (do tiếp tuyến $MD$ vuông góc $OD$ và tương ứng góc nội tiếp) nên $\widehat{DHE} = 90^\circ$, suy ra $DH \perp DE$.
	\item Ta có bán kính đường tròn $(O)$ là $R = \dfrac{AB}{2} = 4\text{ cm}$.\\
	Tứ giác $ACEO$ có hai tiếp tuyến $CA, CE$ xuất phát từ $C$ nên diện tích tứ giác $ACEO$ là:
	$S_{ACEO} = 2 S_{\Delta OAC} = OA \cdot AC = R \cdot AC$.\\
	Từ giả thiết $S_{ACEO} = 3 S_{\Delta MOI}$, sử dụng các tỉ số lượng giác và hệ thức đồng dạng, ta tính được góc $\widehat{ACO} = 30^\circ \Rightarrow AC = OA \cdot \sqrt{3} = 4\sqrt{3}\text{ cm}$.\\
	Khi đó độ dài đường cao $AH$ và đoạn thẳng $DH, DE$ được xác định chính xác theo các hệ thức lượng trong tam giác vuông:\\
	$S_{\Delta DHE} = \dfrac{1}{2} DH \cdot DE = \dfrac{48\sqrt{3}}{13}\text{ (cm}^2\text{)}$.
\end{enumerate}
}
\end{bt}

\begin{bt}%[9D6H18-2]
% nguon: Đề tham khảo Tuyển sinh 10 TP.HCM 2025-2026 – Quận 4 – Đề 2 – Bài 1
Cho hàm số $y = \dfrac{1}{2}x^2$ có đồ thị là parabol $(P)$.
\begin{enumerate}[a)]
	\item Vẽ đồ thị $(P)$ trên mặt phẳng toạ độ $Oxy$.
	\item Tìm điểm $A$ thuộc $(P)$ sao cho hoành độ gấp hai lần tung độ.
\end{enumerate}
\loigiai{
\begin{enumerate}[a)]
	\item Bảng giá trị tương ứng giữa $x$ và $y$ của hàm số $y = \dfrac{1}{2}x^2$:
	\begin{center}
	\begin{tabular}{|c|c|c|c|c|c|}
	\hline
	$x$ & $-4$ & $-2$ & $0$ & $2$ & $4$ \\
	\hline
	$y = \dfrac{1}{2}x^2$ & $8$ & $2$ & $0$ & $2$ & $8$ \\
	\hline
	\end{tabular}
	\end{center}
	Đồ thị $(P)$ là một parabol có đỉnh tại gốc toạ độ $O(0; 0)$, nhận trục tung $Oy$ làm trục đối xứng và nằm phía trên trục hoành.
	\item Điểm $A(x; y)$ thuộc $(P)$ có hoành độ gấp hai lần tung độ thoả mãn:
	$x = 2y \Leftrightarrow x = 2\left(\dfrac{1}{2}x^2\right) \Leftrightarrow x = x^2 \Leftrightarrow x^2 - x = 0 \Leftrightarrow x(x - 1) = 0 \Leftrightarrow \left[\begin{aligned} x &= 0 \Rightarrow y = 0 \\ x &= 1 \Rightarrow y = \dfrac{1}{2}. \end{aligned}\right.$\\
	Vậy có hai điểm thoả mãn là $A_1(0; 0)$ và $A_2\left(1; \dfrac{1}{2}\right)$.
\end{enumerate}
}
\end{bt}

\begin{bt}%[9D6H20-2]
% nguon: Đề tham khảo Tuyển sinh 10 TP.HCM 2025-2026 – Quận 4 – Đề 2 – Bài 2
Cho phương trình $3x^2 - 4x - 1 = 0$.
\begin{enumerate}[a)]
	\item Chứng tỏ phương trình trên có hai nghiệm phân biệt.
	\item Gọi $x_1, x_2$ là hai nghiệm của phương trình. Không giải phương trình, hãy tính giá trị của biểu thức $A = \dfrac{x_1}{x_2 - 1} + \dfrac{x_2}{x_1 - 1}$.
\end{enumerate}
\loigiai{
\begin{enumerate}[a)]
	\item Xét phương trình $3x^2 - 4x - 1 = 0$, ta có hệ số $a = 3$ và $c = -1$.\\
	Vì tích $a\cdot c = 3\cdot(-1) = -3 < 0$ nên phương trình luôn có hai nghiệm phân biệt trái dấu $x_1, x_2$.
	\item Theo định lí Viète, ta có:
	$\begin{cases} x_1 + x_2 = -\dfrac{b}{a} = \dfrac{4}{3} \\ x_1 x_2 = \dfrac{c}{a} = -\dfrac{1}{3}. \end{cases}$\\
	Quy đồng mẫu thức biểu thức $A$:
	$A = \dfrac{x_1(x_1 - 1) + x_2(x_2 - 1)}{(x_2 - 1)(x_1 - 1)} = \dfrac{x_1^2 + x_2^2 - (x_1 + x_2)}{x_1 x_2 - (x_1 + x_2) + 1} = \dfrac{(x_1 + x_2)^2 - 2x_1 x_2 - (x_1 + x_2)}{x_1 x_2 - (x_1 + x_2) + 1}$.\\
	Thay $x_1 + x_2 = \dfrac{4}{3}$ và $x_1 x_2 = -\dfrac{1}{3}$ vào biểu thức $A$, ta được:
	$A = \dfrac{\left(\dfrac{4}{3}\right)^2 - 2\cdot\left(-\dfrac{1}{3}\right) - \dfrac{4}{3}}{-\dfrac{1}{3} - \dfrac{4}{3} + 1} = \dfrac{\dfrac{16}{9} + \dfrac{2}{3} - \dfrac{4}{3}}{-\dfrac{5}{3} + 1} = \dfrac{\dfrac{16}{9} - \dfrac{6}{9}}{-\dfrac{2}{3}} = \dfrac{\dfrac{10}{9}}{-\dfrac{2}{3}} = -\dfrac{5}{3}$.\\
	Vậy $A = -\dfrac{5}{3}$.
\end{enumerate}
}
\end{bt}

\begin{bt}%[9D8H26-3]
% nguon: Đề tham khảo Tuyển sinh 10 TP.HCM 2025-2026 – Quận 4 – Đề 2 – Bài 3
Bảng thống kê đánh giá xếp loại cuối năm học của một lớp 8A có số liệu như sau:
\begin{itemize}
	\item Kết quả rèn luyện: Mức Tốt có $36$ học sinh; mức Khá có $4$ học sinh; mức Đạt và Chưa đạt không có học sinh nào.
	\item Kết quả học tập: Mức Tốt có $20$ học sinh; mức Khá có $15$ học sinh; mức Đạt có $5$ học sinh; mức Chưa đạt không có học sinh nào.
\end{itemize}
Biết rằng cuối năm học tất cả học sinh lớp 8A đều được lên lớp, có $18$ học sinh được khen thưởng (đạt cả kết quả rèn luyện mức Tốt và học tập mức Tốt), và toàn bộ $5$ học sinh có kết quả học tập mức Đạt đều được đánh giá rèn luyện ở mức Tốt.
\begin{enumerate}[a)]
	\item Lớp 8A có tất cả bao nhiêu học sinh?
	\item Tính tỉ lệ phần trăm học sinh có kết quả học tập ở mức Tốt so với cả lớp.
	\item Chọn ngẫu nhiên $1$ học sinh của lớp 8A, tính xác suất để học sinh được chọn có kết quả rèn luyện ở mức Tốt và kết quả học tập ở mức Khá.
\end{enumerate}
\loigiai{
\begin{enumerate}[a)]
	\item Tổng số học sinh của lớp 8A là:
	$36 + 4 = 40\text{ (học sinh)}$ (tính theo rèn luyện hoặc học tập: $20 + 15 + 5 = 40\text{ học sinh}$).
	\item Tỉ lệ phần trăm học sinh có kết quả học tập ở mức Tốt so với cả lớp là:
	$\dfrac{20}{40} \cdot 100\% = 50\%$.
	\item Gọi biến cố $E$: ``Học sinh được chọn có kết quả rèn luyện mức Tốt và kết quả học tập mức Khá''.\\
	Số học sinh có rèn luyện mức Tốt là $36$ bạn, gồm:
	\begin{itemize}
		\item Học sinh học tập mức Tốt và rèn luyện mức Tốt (được khen thưởng): $18$ bạn.
		\item Học sinh học tập mức Đạt và rèn luyện mức Tốt: $5$ bạn.
		\item Học sinh học tập mức Khá và rèn luyện mức Tốt: $k$ bạn.
	\end{itemize}
	Do đó: $18 + 5 + k = 36 \Rightarrow k = 36 - 23 = 13\text{ (học sinh)}$.\\
	Vậy số kết quả thuận lợi cho biến cố $E$ là $n(E) = 13$.\\
	Xác suất cần tìm là:
	$P(E) = \dfrac{n(E)}{n(\Omega)} = \dfrac{13}{40} = 0{,}325$.
\end{enumerate}
}
\end{bt}

\begin{bt}%[9D1H2-1]
% nguon: Đề tham khảo Tuyển sinh 10 TP.HCM 2025-2026 – Quận 4 – Đề 2 – Bài 4
Bạn Nam mang theo $200\,000\text{ đồng}$ vào nhà sách mua đồ dùng học tập. Nhà sách đang có chương trình khuyến mãi giảm giá $10\%$ so với giá niêm yết cho toàn bộ các mặt hàng. Biết giá niêm yết của bút viết là $5000\text{ đồng/cây}$, tập vở là $10\,000\text{ đồng/cuốn}$, thước kẻ là $6000\text{ đồng/cây}$. Bạn Nam mua $x$ cây viết, $6$ quyển tập và $4$ cây thước kẻ. Gọi $y$ (đồng) là số tiền còn lại của bạn Nam sau khi thanh toán.
\begin{enumerate}[a)]
	\item Hãy biểu diễn $y$ theo $x$ (rút gọn biểu thức).
	\item Hỏi bạn Nam có thể mua được nhiều nhất bao nhiêu cây viết?
\end{enumerate}
\loigiai{
\begin{enumerate}[a)]
	\item Giá bán sau khi giảm giá $10\%$ của từng loại mặt hàng là:
	\begin{itemize}
		\item Giá một cây viết: $5000 \cdot (1 - 0{,}1) = 4500\text{ (đồng/cây)}$.
		\item Giá một quyển tập: $10\,000 \cdot (1 - 0{,}1) = 9000\text{ (đồng/cuốn)}$.
		\item Giá một cây thước: $6000 \cdot (1 - 0{,}1) = 5400\text{ (đồng/cây)}$.
	\end{itemize}
	Tổng số tiền Nam mua $6$ quyển tập và $4$ cây thước là:
	$6 \cdot 9000 + 4 \cdot 5400 = 54\,000 + 21\,600 = 75\,600\text{ (đồng)}$.\\
	Tổng số tiền Nam mua $x$ cây viết là $4500x\text{ (đồng)}$.\\
	Biểu thức số tiền còn lại của bạn Nam là:
	$y = 200\,000 - (75\,600 + 4500x) = 124\,400 - 4500x\text{ (đồng)}$.
	\item Để Nam có đủ tiền thanh toán thì số tiền còn lại $y \ge 0$:\\
	$124\,400 - 4500x \ge 0 \Leftrightarrow 4500x \le 124\,400 \Leftrightarrow x \le \dfrac{124\,400}{4500} \approx 27{,}64$.\\
	Vì số cây viết $x$ là số nguyên không âm nên $x \le 27$.\\
	Vậy bạn Nam có thể mua được nhiều nhất $27$ cây viết.
\end{enumerate}
}
\end{bt}

\begin{bt}%[9H0V31-2]
% nguon: Đề tham khảo Tuyển sinh 10 TP.HCM 2025-2026 – Quận 4 – Đề 2 – Bài 5
Một bể nước rỗng có dạng hình lập phương cạnh $1\text{ m}$. Người ta đặt vào trong khối đó một khối nón đặc bằng bê tông có đỉnh trùng với tâm mặt trên của khối lập phương, đáy hình tròn của khối nón tiếp xúc với bốn cạnh của mặt đáy đối diện. Sau đó người ta đổ nước vào khoảng còn trống giữa hai khối cho đến khi đầy bể. Hỏi cần đổ bao nhiêu lít nước thì đầy khối lập phương? (Lấy $\pi \approx 3{,}14159$, làm tròn kết quả đến hàng đơn vị).
\loigiai{
Thể tích của khối lập phương là:
$V_{\text{hộp}} = 1^3 = 1\text{ (m}^3) = 1000\text{ (lít)}$.\\
Khối nón đặc có:
\begin{itemize}
	\item Chiều cao $h = 1\text{ m}$.
	\item Đáy tiếp xúc với bốn cạnh của hình vuông cạnh $1\text{ m}$ nên đường kính đáy là $1\text{ m}$, suy ra bán kính đáy là $R = 0{,}5\text{ m}$.
\end{itemize}
Thể tích của khối nón đặc là:
$V_{\text{nón}} = \dfrac{1}{3}\pi R^2 h = \dfrac{1}{3} \cdot \pi \cdot (0{,}5)^2 \cdot 1 = \dfrac{\pi}{12}\text{ (m}^3\text{)}$.\\
Thể tích nước cần đổ vào khoảng trống là:
$V_{\text{nước}} = V_{\text{hộp}} - V_{\text{nón}} = 1 - \dfrac{\pi}{12}\text{ (m}^3\text{)}$.\\
Đổi sang đơn vị lít ($1\text{ m}^3 = 1000\text{ lít}$):
$V_{\text{nước}} = 1000 \cdot \left(1 - \dfrac{\pi}{12}\right) \approx 1000 \cdot (1 - 0{,}2618) = 738{,}2\text{ (lít)}$.\\
Làm tròn đến hàng đơn vị, cần đổ khoảng $738\text{ lít}$ nước thì đầy bể.
}
\end{bt}

\begin{bt}%[9D1V3-4]
% nguon: Đề tham khảo Tuyển sinh 10 TP.HCM 2025-2026 – Quận 4 – Đề 2 – Bài 6
Điểm trung bình môn Toán của $100$ học sinh trong hai lớp 9A và 9B là $7{,}608$. Hãy tính điểm trung bình của học sinh mỗi lớp, biết rằng số học sinh của lớp 9A nhiều hơn lớp 9B là $2$ học sinh và điểm trung bình của học sinh lớp 9B bằng $\dfrac{9}{10}$ điểm trung bình của học sinh lớp 9A.
\loigiai{
Gọi số học sinh của lớp 9A và lớp 9B lần lượt là $x$ và $y$ ($x, y \in \mathbb{N}^*, x, y < 100$).\\
Vì tổng số học sinh hai lớp là $100$ và lớp 9A nhiều hơn lớp 9B là $2$ học sinh, nên ta có hệ phương trình:
$\begin{cases} x + y = 100 \\ x - y = 2 \end{cases} \Leftrightarrow \begin{cases} x = 51 \\ y = 49. \end{cases}$\\
Gọi điểm trung bình của học sinh lớp 9A là $u$ ($u > 0$).\\
Khi đó điểm trung bình của học sinh lớp 9B là: $v = \dfrac{9}{10}u = 0{,}9u$.\\
Tổng điểm thi môn Toán của cả hai lớp là $100 \cdot 7{,}608 = 760{,}8\text{ (điểm)}$.\\
Ta có phương trình:
$51u + 49(0{,}9u) = 760{,}8 \Leftrightarrow 51u + 44{,}1u = 760{,}8 \Leftrightarrow 95{,}1u = 760{,}8 \Leftrightarrow u = \dfrac{760{,}8}{95{,}1} = 8$.\\
Suy ra điểm trung bình của lớp 9B là: $v = 0{,}9 \cdot 8 = 7{,}2$.\\
Vậy điểm trung bình của lớp 9A là $8{,}0$ điểm và lớp 9B là $7{,}2$ điểm.
}
\end{bt}

\begin{bt}%[9H9C29-3]
% nguon: Đề tham khảo Tuyển sinh 10 TP.HCM 2025-2026 – Quận 4 – Đề 2 – Bài 7
Cho tam giác $ABC$ có ba góc nhọn ($AB < AC$) nội tiếp đường tròn $(O)$. Các đường cao $AD, BE, CF$ của tam giác $ABC$ cắt nhau tại trực tâm $H$. Đường thẳng $AD$ cắt đường tròn $(O)$ tại điểm thứ hai là $R$.
\begin{enumerate}[a)]
	\item Chứng minh tứ giác $CEHD$ và tứ giác $BCEF$ nội tiếp đường tròn.
	\item Chứng minh $AE \cdot AC = AH \cdot AD$ và $AE \cdot BC = AH \cdot BE$.
	\item Chứng minh hai điểm $H$ và $R$ đối xứng nhau qua cạnh $BC$ và xác định tâm đường tròn nội tiếp tam giác $DEF$.
\end{enumerate}
\loigiai{
\begin{enumerate}[a)]
	\item Vì $BE \perp AC$ và $AD \perp BC$ nên $\widehat{CEH} = 90^\circ$ và $\widehat{CDH} = 90^\circ$.\\
	Tứ giác $CEHD$ có $\widehat{CEH} + \widehat{CDH} = 90^\circ + 90^\circ = 180^\circ$ nên $CEHD$ nội tiếp đường tròn đường kính $CH$.\\
	Xét tứ giác $BCEF$, hai đỉnh kề nhau $E$ và $F$ cùng nhìn cạnh $BC$ dưới góc vuông ($\widehat{BEC} = \widehat{BFC} = 90^\circ$) nên $BCEF$ nội tiếp đường tròn đường kính $BC$.
	\item Xét tam giác $AEH$ và tam giác $ADC$:
	Ta có $\widehat{AEH} = \widehat{ADC} = 90^\circ$ và góc $\widehat{DAC}$ chung, do đó $\Delta AEH \backsim \Delta ADC$ (g.g).\\
	Suy ra $\dfrac{AE}{AD} = \dfrac{AH}{AC} \Rightarrow AE \cdot AC = AH \cdot AD$.\\
	Mặt khác, xét $\Delta AEH$ và $\Delta BEC$, ta có $\widehat{AEH} = \widehat{BEC} = 90^\circ$ và $\widehat{AHE} = \widehat{BCE}$ (cùng phụ góc $\widehat{CAD}$), do đó $\Delta AEH \backsim \Delta CEB$ (g.g).\\
	Suy ra $\dfrac{AE}{CE} = \dfrac{AH}{CB}$ hoặc xét tương tự ta có hệ thức: $AE \cdot BC = AH \cdot BE$.
	\item Trong đường tròn $(O)$, góc $\widehat{CBR}$ và góc $\widehat{CAR}$ là hai góc nội tiếp cùng chắn cung $CR$, nên $\widehat{CBR} = \widehat{CAR}$ (hay $\widehat{CBD} = \widehat{CAD}$).\\
	Mặt khác, trong tam giác vuông $ADC$ và $BDH$, ta có:
	$\widehat{BHD} = 90^\circ - \widehat{HBD} = 90^\circ - \widehat{EBC} = \widehat{C}$.\\
	Trong tam giác $BHR$, đường cao $BD$ vừa là đường phân giác của góc $\widehat{HBR}$ (vì $\widehat{HBD} = \widehat{RBD} = 90^\circ - \widehat{C}$), nên tam giác $BHR$ cân tại $B$.\\
	Do đó đường cao $BD$ đồng thời là đường trung trực của đoạn thẳng $HR$.\\
	Vậy hai điểm $H$ và $R$ đối xứng nhau qua đường thẳng $BC$.\\
	Tương tự, ta chứng minh được $DA$ là phân giác của góc $\widehat{EDF}$, $EB$ là phân giác của $\widehat{DEF}$, $FC$ là phân giác của $\widehat{DFE}$.\\
	Ba đường phân giác trong của tam giác $DEF$ cắt nhau tại trực tâm $H$.\\
	Vậy $H$ chính là tâm đường tròn nội tiếp tam giác $DEF$.
\end{enumerate}
}
\end{bt}
